% Problemas del capítulo 11. Generado por tools/problemas/gen.py; no editar a mano.
\section{Problemas}\label{sec:problemas-11}

Graduados: \(\star\) directo, \(\star\star\) para pensar, \(\star\star\star\) desafío. Las respuestas están en el apéndice~\ref{ap:respuestas}; conviene intentarlos antes de mirarlas.

\begin{problema}[\({\star}\)]\label{pb:11-1}
Calculá la inversa de \(A=\begin{pmatrix}2&1\\5&3\end{pmatrix}\) y comprobá que \(AA^{-1}=I\).
\end{problema}

\begin{problema}[\({\star}{\star}\)]\label{pb:11-2}
Sea \(R\) el giro de \(90^\circ\) y \(S=\mathrm{diag}(2,1)\). La transformación «primero \(S\), después \(R\)» es \(RS\). ¿Cuál es su inversa? Comprobalo con el punto \((1,1)\).
\end{problema}

\begin{problema}[\({\star}{\star}\)]\label{pb:11-3}
Buscá la recta \(y=a+bx\) que mejor pasa por \((0,1)\), \((1,2)\) y \((2,2)\), por cuadrados mínimos verticales. Calculá los residuos y comprobá que suman cero.
\end{problema}

\begin{problema}[\({\star}{\star}\)]\label{pb:11-4}
Se conoce la altura del punto \(A\): \(10{,}000\)~m. Se mide cuánto más alto está cada punto que el anterior (el desnivel) en un recorrido que vuelve al punto de partida: \(A\to B=+1{,}250\), \(B\to C=+0{,}830\) y \(C\to A=-2{,}074\)~m. Si las medidas fueran perfectas, sumarían cero. ¿Por cuánto no cierran? Repartí la diferencia en partes iguales y calculá las alturas de \(B\) y \(C\).
\end{problema}

\begin{problema}[\({\star}{\star}{\star}\)]\label{pb:11-5}
\(A=\begin{pmatrix}1&1\\1&1\end{pmatrix}\) no tiene inversa. Calculá su pseudoinversa \(A^+\) y la solución \(x^+=A^+b\) para \(b=(2,4)\). ¿Qué es \(Ax^+\)?
\end{problema}
